\section{坐标系、坐标轴上的单位矢量与矢量基}

\subsection{右手笛卡尔直角坐标系及其坐标轴上的单位矢量}

% 图：右手笛卡尔直角坐标系，坐标轴上的单位矢量 $\vec{e}_1, \vec{e}_2, \vec{e}_3$

右手笛卡尔直角坐标系坐标轴上的单位矢量满足
\begin{equation}
\begin{cases}
\vec{e}_1 \cdot \vec{e}_1 = 1; \quad \vec{e}_1 \cdot \vec{e}_2 = 0; \quad \vec{e}_1 \cdot \vec{e}_3 = 0; \\
\vec{e}_2 \cdot \vec{e}_1 = 0; \quad \vec{e}_2 \cdot \vec{e}_2 = 1; \quad \vec{e}_2 \cdot \vec{e}_3 = 0; \\
\vec{e}_3 \cdot \vec{e}_1 = 0; \quad \vec{e}_3 \cdot \vec{e}_2 = 0; \quad \vec{e}_3 \cdot \vec{e}_3 = 1.
\end{cases}
\end{equation}

即
\begin{equation}\label{eq:4.kronecker}
\vec{e}_m \cdot \vec{e}_n = \delta_{mn}
\end{equation}

其中 Kronecker $\delta$ 符号定义为
\begin{equation}\label{eq:4.delta_def}
\delta_{mn} = \begin{cases} 1, & m = n \\ 0, & m \neq n \end{cases}
\end{equation}

右手笛卡尔直角坐标系坐标轴上的单位矢量的叉积满足
\begin{equation}
\begin{cases}
\vec{e}_1 \times \vec{e}_1 = \vec{0}; \quad \vec{e}_1 \times \vec{e}_2 = \vec{e}_3; \quad \vec{e}_1 \times \vec{e}_3 = -\vec{e}_2; \\
\vec{e}_2 \times \vec{e}_1 = -\vec{e}_3; \quad \vec{e}_2 \times \vec{e}_2 = \vec{0}; \quad \vec{e}_2 \times \vec{e}_3 = \vec{e}_1; \\
\vec{e}_3 \times \vec{e}_1 = \vec{e}_2; \quad \vec{e}_3 \times \vec{e}_2 = -\vec{e}_1; \quad \vec{e}_3 \times \vec{e}_3 = \vec{0}.
\end{cases}
\end{equation}

即
\begin{equation}\label{eq:4.levi_civita}
\vec{e}_m \times \vec{e}_n = \sum_{l=1}^{3} \varepsilon_{mnl}\, \vec{e}_l =:\, \varepsilon_{mnl}\, \vec{e}_l
\end{equation}

其中 Levi-Civita 符号（排列符号）定义为
\begin{equation}\label{eq:4.epsilon_def}
\varepsilon_{mnl} = \begin{cases}
1, & (m, n, l) = (1, 2, 3) \text{ 或 } (2, 3, 1) \text{ 或 } (3, 1, 2) \\
-1, & (m, n, l) = (2, 1, 3) \text{ 或 } (3, 2, 1) \text{ 或 } (1, 3, 2) \\
0, & m = n \text{ 或 } n = l \text{ 或 } l = m
\end{cases}
\end{equation}


\subsection{右手笛卡尔直角坐标系矢量基及其运算法则}

% 图：右手笛卡尔直角坐标系矢量基

右手笛卡尔直角坐标系矢量基
\begin{equation}
\boldsymbol{e} \triangleq \begin{bmatrix} \vec{e}_1 \\ \vec{e}_2 \\ \vec{e}_3 \end{bmatrix}, \qquad \boldsymbol{e}^{\mathrm{T}} = [\vec{e}_1,\, \vec{e}_2,\, \vec{e}_3]
\end{equation}

坐标系矢量基运算法则：
\begin{equation}
\boldsymbol{e}^{\mathrm{T}} \cdot \boldsymbol{e} = [\vec{e}_1,\, \vec{e}_2,\, \vec{e}_3] \cdot \begin{bmatrix} \vec{e}_1 \\ \vec{e}_2 \\ \vec{e}_3 \end{bmatrix} = \vec{e}_1 \cdot \vec{e}_1 + \vec{e}_2 \cdot \vec{e}_2 + \vec{e}_3 \cdot \vec{e}_3 = 3
\end{equation}

\begin{equation}
\boldsymbol{e} \cdot \boldsymbol{e}^{\mathrm{T}} = \begin{bmatrix} \vec{e}_1 \\ \vec{e}_2 \\ \vec{e}_3 \end{bmatrix} \cdot [\vec{e}_1,\, \vec{e}_2,\, \vec{e}_3] = \begin{bmatrix} \vec{e}_1 \cdot \vec{e}_1 & \vec{e}_1 \cdot \vec{e}_2 & \vec{e}_1 \cdot \vec{e}_3 \\ \vec{e}_2 \cdot \vec{e}_1 & \vec{e}_2 \cdot \vec{e}_2 & \vec{e}_2 \cdot \vec{e}_3 \\ \vec{e}_3 \cdot \vec{e}_1 & \vec{e}_3 \cdot \vec{e}_2 & \vec{e}_3 \cdot \vec{e}_3 \end{bmatrix} = \boldsymbol{I}_3
\end{equation}

\begin{equation}
\boldsymbol{e}^{\mathrm{T}} \times \boldsymbol{e} = [\vec{e}_1,\, \vec{e}_2,\, \vec{e}_3] \times \begin{bmatrix} \vec{e}_1 \\ \vec{e}_2 \\ \vec{e}_3 \end{bmatrix} = \vec{e}_1 \times \vec{e}_1 + \vec{e}_2 \times \vec{e}_2 + \vec{e}_3 \times \vec{e}_3 = \vec{0}
\end{equation}

\begin{equation}
\boldsymbol{e} \times \boldsymbol{e}^{\mathrm{T}} = \begin{bmatrix} \vec{e}_1 \\ \vec{e}_2 \\ \vec{e}_3 \end{bmatrix} \times [\vec{e}_1,\, \vec{e}_2,\, \vec{e}_3] = \begin{bmatrix} \vec{e}_1 \times \vec{e}_1 & \vec{e}_1 \times \vec{e}_2 & \vec{e}_1 \times \vec{e}_3 \\ \vec{e}_2 \times \vec{e}_1 & \vec{e}_2 \times \vec{e}_2 & \vec{e}_2 \times \vec{e}_3 \\ \vec{e}_3 \times \vec{e}_1 & \vec{e}_3 \times \vec{e}_2 & \vec{e}_3 \times \vec{e}_3 \end{bmatrix} = \begin{bmatrix} \vec{0} & \vec{e}_3 & -\vec{e}_2 \\ -\vec{e}_3 & \vec{0} & \vec{e}_1 \\ \vec{e}_2 & -\vec{e}_1 & \vec{0} \end{bmatrix}
\end{equation}

对应任意三维列阵 $\boldsymbol{u} = [u_1,\, u_2,\, u_3]^{\mathrm{T}}$，引入如下定义式
\begin{equation}\label{eq:4.skew_matrix}
\tilde{\boldsymbol{u}} \triangleq \begin{bmatrix} 0 & -u_3 & u_2 \\ u_3 & 0 & -u_1 \\ -u_2 & u_1 & 0 \end{bmatrix}, \qquad \tilde{u}(m, n) \triangleq -\varepsilon_{mnl}\, u_l
\end{equation}

满足
\begin{align}
\tilde{\boldsymbol{u}}^{\mathrm{T}} &= -\tilde{\boldsymbol{u}} \\
\tilde{\boldsymbol{u}}\, \boldsymbol{v} &= -\tilde{\boldsymbol{v}}\, \boldsymbol{u}, \quad \boldsymbol{v}^{\mathrm{T}} \tilde{\boldsymbol{u}} = -\boldsymbol{u}^{\mathrm{T}} \tilde{\boldsymbol{v}} \\
\tilde{\boldsymbol{u}}\, \boldsymbol{u} &= \boldsymbol{0} \\
\tilde{\boldsymbol{u}}\, \tilde{\boldsymbol{v}} &= \boldsymbol{v}\, \boldsymbol{u}^{\mathrm{T}} - (\boldsymbol{v}^{\mathrm{T}} \boldsymbol{u})\, \boldsymbol{I}_3
\end{align}

由此可得
\begin{equation}
\boldsymbol{e} \times \boldsymbol{e}^{\mathrm{T}} = \begin{bmatrix} \vec{0} & \vec{e}_3 & -\vec{e}_2 \\ -\vec{e}_3 & \vec{0} & \vec{e}_1 \\ \vec{e}_2 & -\vec{e}_1 & \vec{0} \end{bmatrix} = -\tilde{\boldsymbol{e}}
\end{equation}

汇总：
\begin{myformula}
\begin{align}
\boldsymbol{e}^{\mathrm{T}} \cdot \boldsymbol{e} &= 3 \tag{I} \\
\boldsymbol{e} \cdot \boldsymbol{e}^{\mathrm{T}} &= \boldsymbol{I}_3 \tag{II} \\
\boldsymbol{e}^{\mathrm{T}} \times \boldsymbol{e} &= \vec{0} \tag{III} \\
\boldsymbol{e} \times \boldsymbol{e}^{\mathrm{T}} &= -\tilde{\boldsymbol{e}} \tag{IV}
\end{align}
\end{myformula}


\section{矢量及矢量运算的坐标表示}

\subsection{矢量的坐标表示}

矢量的坐标分量表示形式
\begin{equation}
\vec{u} = u_1\, \vec{e}_1 + u_2\, \vec{e}_2 + u_3\, \vec{e}_3
\end{equation}

可记为矩阵乘积（数乘）的形式，即
\begin{equation}\label{eq:4.vector_rep}
\vec{u} = [\vec{e}_1,\, \vec{e}_2,\, \vec{e}_3] \begin{bmatrix} u_1 \\ u_2 \\ u_3 \end{bmatrix} =:\, \boldsymbol{e}^{\mathrm{T}} \boldsymbol{u} = \boldsymbol{u}^{\mathrm{T}} \boldsymbol{e}
\end{equation}

式中
\begin{equation}
\boldsymbol{u} = \begin{bmatrix} u_1 \\ u_2 \\ u_3 \end{bmatrix}
\end{equation}

称为矢量 $\vec{u}$ 的坐标列阵或投影列阵。

% 图：矢量 $\vec{u}$ 在坐标系中的表示


\subsection{两矢量点积运算的坐标表示}

两矢量的点积（内积）：
\begin{equation}\label{eq:4.dot_product}
\vec{u} \cdot \vec{v} = \boldsymbol{u}^{\mathrm{T}} \boldsymbol{e} \cdot \boldsymbol{e}^{\mathrm{T}} \boldsymbol{v} = \boldsymbol{u}^{\mathrm{T}} \boldsymbol{v}
\end{equation}

矢量点积的性质：
\begin{equation}
\vec{u} \cdot \vec{u} = \boldsymbol{u}^{\mathrm{T}} \boldsymbol{u} = |\vec{u}|^2 \tag{I}
\end{equation}

\begin{equation}
\vec{u} \cdot \vec{v} = \boldsymbol{u}^{\mathrm{T}} \boldsymbol{v} = \boldsymbol{v}^{\mathrm{T}} \boldsymbol{u} = \vec{v} \cdot \vec{u} \tag{II}
\end{equation}

\begin{equation}
(\vec{u} - \vec{v}) \cdot (\vec{u} - \vec{v}) = \vec{u} \cdot \vec{u} - 2\, \vec{u} \cdot \vec{v} + \vec{v} \cdot \vec{v} = |\vec{u}|^2 + |\vec{v}|^2 - 2\, \vec{u} \cdot \vec{v} = |\vec{u} - \vec{v}|^2 = |\vec{u}|^2 + |\vec{v}|^2 - 2\, |\vec{u}|\, |\vec{v}| \cos\langle \vec{u},\, \vec{v} \rangle
\end{equation}

\begin{equation}
\vec{u} \cdot \vec{v} = |\vec{u}|\, |\vec{v}| \cos\langle \vec{u},\, \vec{v} \rangle \tag{III}
\end{equation}


\subsection{两矢量叉积运算的坐标表示}

两矢量的叉积：
\begin{equation}\label{eq:4.cross_product}
\vec{u} \times \vec{v} = \boldsymbol{u}^{\mathrm{T}} \boldsymbol{e} \times \boldsymbol{e}^{\mathrm{T}} \boldsymbol{v} = -\boldsymbol{u}^{\mathrm{T}} \tilde{\boldsymbol{e}}\, \boldsymbol{v} = \boldsymbol{e}^{\mathrm{T}} \tilde{\boldsymbol{u}}\, \boldsymbol{v} =:\, \boldsymbol{e}^{\mathrm{T}} [\boldsymbol{u}]_\times \boldsymbol{v}
\end{equation}

矢量叉乘的性质：
\begin{equation}
\vec{u} \times \vec{v} = \boldsymbol{e}^{\mathrm{T}} \tilde{\boldsymbol{u}}\, \boldsymbol{v} = -\boldsymbol{e}^{\mathrm{T}} \tilde{\boldsymbol{v}}\, \boldsymbol{u} = -\vec{v} \times \vec{u} \tag{I}
\end{equation}

\begin{equation}
\vec{u} \cdot (\vec{u} \times \vec{v}) = \boldsymbol{u}^{\mathrm{T}} \tilde{\boldsymbol{u}}\, \boldsymbol{v} = 0 \tag{II}
\end{equation}

\begin{equation}
\vec{v} \cdot (\vec{u} \times \vec{v}) = -\boldsymbol{v}^{\mathrm{T}} \tilde{\boldsymbol{v}}\, \boldsymbol{u} = 0 \tag{III}
\end{equation}

\begin{align}
|\vec{u} \times \vec{v}|^2 &= (\vec{u} \times \vec{v}) \cdot (\vec{u} \times \vec{v}) = \boldsymbol{v}^{\mathrm{T}} \tilde{\boldsymbol{u}}^{\mathrm{T}} \tilde{\boldsymbol{u}}\, \boldsymbol{v} = |\vec{u}|^2\, |\vec{v}|^2 - (\vec{u} \cdot \vec{v})^2 \nonumber \\
&= |\vec{u}|^2\, |\vec{v}|^2 - |\vec{u}|^2\, |\vec{v}|^2 \cos^2\langle \vec{u},\, \vec{v} \rangle = |\vec{u}|^2\, |\vec{v}|^2 \sin^2\langle \vec{u},\, \vec{v} \rangle
\end{align}

\begin{equation}
|\vec{u} \times \vec{v}| = |\vec{u}|\, |\vec{v}| \sin\langle \vec{u},\, \vec{v} \rangle \tag{IV}
\end{equation}


\subsection{三矢量混合积运算的坐标表示}

三矢量的混合积：
\begin{equation}\label{eq:4.triple_product}
\vec{u} \cdot (\vec{v} \times \vec{w}) = \boldsymbol{u}^{\mathrm{T}} \boldsymbol{e} \cdot (\boldsymbol{e}^{\mathrm{T}} \tilde{\boldsymbol{v}}\, \boldsymbol{w}) = \boldsymbol{u}^{\mathrm{T}} \boldsymbol{e} \cdot \boldsymbol{e}^{\mathrm{T}} \tilde{\boldsymbol{v}}\, \boldsymbol{w} = \boldsymbol{u}^{\mathrm{T}} \tilde{\boldsymbol{v}}\, \boldsymbol{w}
\end{equation}

三矢量混合积的性质：
\begin{align}
\vec{u} \cdot (\vec{v} \times \vec{w}) &= \boldsymbol{u}^{\mathrm{T}} \tilde{\boldsymbol{v}}\, \boldsymbol{w} = \boldsymbol{v}^{\mathrm{T}} \tilde{\boldsymbol{w}}\, \boldsymbol{u} = \vec{v} \cdot (\vec{w} \times \vec{u}) \nonumber \\
&= \boldsymbol{w}^{\mathrm{T}} \tilde{\boldsymbol{u}}\, \boldsymbol{v} = \vec{w} \cdot (\vec{u} \times \vec{v})
\end{align}


\subsection{三矢量叉积运算的坐标表示}

三矢量的叉积：
\begin{equation}\label{eq:4.double_cross}
\vec{u} \times (\vec{v} \times \vec{w}) = (\boldsymbol{e}^{\mathrm{T}} \boldsymbol{u}) \times (\boldsymbol{e}^{\mathrm{T}} \tilde{\boldsymbol{v}}\, \boldsymbol{w}) = \boldsymbol{e}^{\mathrm{T}} \tilde{\boldsymbol{u}}\, \tilde{\boldsymbol{v}}\, \boldsymbol{w}
\end{equation}

三矢量叉积的性质：
\begin{align}
\vec{u} \times (\vec{v} \times \vec{w}) &= \boldsymbol{e}^{\mathrm{T}} \tilde{\boldsymbol{u}}\, \tilde{\boldsymbol{v}}\, \boldsymbol{w} = \boldsymbol{e}^{\mathrm{T}} \left[\boldsymbol{v}\, \boldsymbol{u}^{\mathrm{T}} - (\boldsymbol{v}^{\mathrm{T}} \boldsymbol{u})\, \boldsymbol{I}_3\right] \boldsymbol{w} = \boldsymbol{e}^{\mathrm{T}} \left[\boldsymbol{v}\, \boldsymbol{u}^{\mathrm{T}} \boldsymbol{w} - (\boldsymbol{v}^{\mathrm{T}} \boldsymbol{u})\, \boldsymbol{w}\right] \nonumber \\
&= \boldsymbol{e}^{\mathrm{T}} \left[\boldsymbol{v}\, (\boldsymbol{u}^{\mathrm{T}} \boldsymbol{w}) - \boldsymbol{w}\, (\boldsymbol{v}^{\mathrm{T}} \boldsymbol{u})\right] = \boldsymbol{e}^{\mathrm{T}} \boldsymbol{v}\, (\boldsymbol{u}^{\mathrm{T}} \boldsymbol{w}) - \boldsymbol{e}^{\mathrm{T}} \boldsymbol{w}\, (\boldsymbol{v}^{\mathrm{T}} \boldsymbol{u}) \nonumber \\
&= \vec{v}\, (\boldsymbol{u}^{\mathrm{T}} \boldsymbol{w}) - \vec{w}\, (\boldsymbol{v}^{\mathrm{T}} \boldsymbol{u}) = (\vec{u} \cdot \vec{w})\, \vec{v} - (\vec{u} \cdot \vec{v})\, \vec{w}
\end{align}
